\section{Ordinal operations and the well-ordering threshold}
\label{sec:p03-ordinal-pinning}

Work with the one-sorted regular logics, nonempty structures and
arbitrary set-sized vocabularies of
\cref{def:p03-general-pinning-setting}. Let \(\kappa\geq1\)
be a cardinal. A theory has at most \(\kappa\) sentences; the
vocabulary is not subject to this bound.

\begin{hbdefinition}{def:p03-op-carrier}{Carrier pinning}
A theory \(T\), with a designated unary predicate \(P\) and binary
relation \(<\), \emph{pins an ordinal on a carrier} if every model
of \(T\) makes \(<\) restricted to \(P\) a strict well-order,
and some model gives that carrier the specified order type.
The carrier may be empty even though the ambient structure is not.
Write \(W_\kappa\) for the supremum of these ordinals, or
\(\infty\) when they are unbounded in the ordinals.
\end{hbdefinition}

\begin{hbtheorem}{thm:p03-op-initial}{proved=annex:p03.op.1}{Initial segments, successors and carriers}
Carrier-pinnable ordinals are closed downwards and under successor.
Every finite ordinal is carrier-pinnable with one sentence.
Moreover \(W_\kappa=w_\kappa(\mathscr L)\), where the latter
uses the field convention of \cref{def:p03-w03-p03-041-pinning}.
If this common value is an ordinal \(W\), it is an infinite limit,
is not itself pinnable, and the carrier-pinnable ordinals are
exactly those below \(W\). Under the strict-field convention,
exactly the same ordinals are pinnable except for \(1\).
\end{hbtheorem}

\begin{hbtheorem}{thm:p03-op-threshold}{proved=annex:p03.op.2}{The cofinal-model threshold}
If \(W_\kappa=W\) is an ordinal, then \(W\) is the least ordinal
\(\gamma\) with the following property: whenever a theory of at
most \(\kappa\) sentences has well-ordered carrier models whose
order types below \(\gamma\) are cofinal in \(\gamma\), it also
has a model whose designated carrier order is not well-ordered.
If \(W_\kappa=\infty\), no ordinal has this property.
For strict fields rather than designated carriers, \(W\) is the
least \emph{infinite} ordinal with the corresponding property.
\end{hbtheorem}

The qualification about fields is necessary: the field of a strict
order cannot have exactly one element. Consequently the condition
on types cofinal below \(2\) is impossible for fields, and the
unqualified least-ordinal characterization would be vacuous there.
Carrier pinning gives the full threshold formulation without
changing the well-ordering number. Compare
\cite[Ch.~II, Proposition~5.2.2, pp.~61--62]{Ebbinghaus1985ExtendedLogics}.

\begin{hbtheorem}{thm:p03-op-arithmetic}{proved=annex:p03.op.3}{Ordinal sum, product and exponentiation}
If \(\alpha\) and \(\beta\) are carrier-pinnable by theories of
at most \(\kappa\) sentences, so are
\(\alpha+\beta\), \(\alpha\beta\), and \(\alpha^\beta\),
with \(0^0=1\). In particular, a bounded \(W_\kappa\) is
closed under these operations on its smaller ordinals.
\end{hbtheorem}
